\chapter{Charles Bonnet Syndrome}

\section*{Definition and Overview}
Charles Bonnet syndrome is a condition in which people who have lost part or most of their sight experience vivid, complex visual hallucinations. These hallucinations may take the form of patterns, shapes, faces, people, animals, buildings, or entire scenes, and the person usually recognises that what they are seeing is not real. It is not a mental illness and is not a sign of dementia. The syndrome is named after the Swiss naturalist who first described it in his grandfather in the eighteenth century. It is considerably more common than is generally appreciated, and many affected people never mention their symptoms for fear of being thought mentally unwell.

\section*{Epidemiology}
Charles Bonnet syndrome occurs in people with significant vision loss from any cause, and its frequency rises with the severity of sight loss. Studies suggest that approximately one in eight to one in ten people with severe visual impairment experience visual hallucinations at some point, although the true figure may be higher because many never report them. It is most common in older adults, simply because age-related eye disease is so common, but it can occur at any age, including in children. Typical underlying causes include age-related macular degeneration, glaucoma, diabetic retinopathy, cataracts, and damage to the visual pathways from stroke or other neurological injury.

\section*{Aetiology and Pathophysiology}
\medskip
\textbf{Mechanism.} Charles Bonnet syndrome arises in the brain of a person with poor vision, rather than in the eyes themselves:
\begin{itemize}
    \item \textbf{Sensory deprivation:} when the brain receives markedly reduced visual input, its higher processing areas become overactive and "fill in" missing information with images
    \item \textbf{Release phenomenon:} without normal input, previously suppressed regions of the visual cortex generate spontaneous activity that the person experiences as hallucinations
    \item \textbf{Underlying causes:} any significant loss of vision, including macular degeneration, glaucoma, cataracts, diabetic retinopathy, and stroke affecting the visual pathway
    \item \textbf{Not psychiatric:} the hallucinations occur in people with otherwise intact thinking, memory, and insight
    \item \textbf{Worsening factors:} dim light, tiredness, stress, and social isolation often make the hallucinations more frequent or vivid
\end{itemize}

\section*{Clinical Features}
\begin{itemize}
    \item \textbf{Content:} hallucinations range from simple patterns and geometric shapes to detailed faces, people, animals, buildings, or whole scenes
    \item \textbf{Timing:} images appear suddenly and last from seconds to several hours
    \item \textbf{Insight:} the person usually knows that the images are not real
    \item \textbf{Appearance:} images may be in colour or black and white, still or moving, and may enlarge or repeat
    \item \textbf{Tone:} most hallucinations are neutral or pleasant, although some people find them disturbing
    \item \textbf{Modality:} the hallucinations are purely visual -- there are no voices, smells, or other false perceptions
    \item \textbf{Aggravating situations:} symptoms are worse in dim light, in quiet settings, or when alone, and often ease with brighter lighting or activity
    \item \textbf{Cognition:} there is no confusion, memory loss, or disordered thinking accompanying the hallucinations themselves
\end{itemize}

\section*{Differential Diagnosis}
\begin{itemize}
    \item Delirium -- sudden confusion with hallucinations in a medically unwell person
    \item Dementia, particularly dementia with Lewy bodies
    \item Schizophrenia and other psychotic disorders
    \item Drug-induced hallucinations from medicines, alcohol, or recreational drugs
    \item Migraine with visual aura
    \item Focal epilepsy affecting the visual areas of the brain
    \item Brain tumours or stroke affecting the visual cortex
    \item Parkinson's disease and the medicines used to treat it
\end{itemize}

\section*{When to Seek Medical Care}
See a doctor if you experience visual hallucinations for the first time, so that other causes such as confusion, medication effects, or neurological disease can be excluded. People who already have a diagnosis of Charles Bonnet syndrome should still be reviewed if the hallucinations become more frequent or distressing, or if they are accompanied by other new symptoms.

\medskip
\textbf{Call 000 (or local emergency services) for:}
\begin{itemize}
    \item Sudden confusion, drowsiness, or a change in mental state
    \item Hallucinations accompanied by a severe headache, weakness, or difficulty speaking
    \item Sudden complete or severe loss of vision
    \item Hallucinations accompanied by a seizure or fit
\end{itemize}

\section*{Diagnosis and Investigations}
\begin{itemize}
    \item \textbf{Clinical history:} the diagnosis is made chiefly by discussing the hallucinations -- their content, timing, frequency, and the person's insight into them
    \item \textbf{Eye examination:} assessment of the degree and cause of vision loss by an optometrist or ophthalmologist
    \item \textbf{Cognitive assessment:} checking that memory, thinking, and orientation are intact, to exclude dementia or delirium
    \item \textbf{Medication review:} to identify drugs that can provoke hallucinations
    \item \textbf{Mental health assessment:} to rule out psychiatric causes in uncertain cases
    \item \textbf{Investigations:} blood tests and brain imaging are only needed if another neurological or medical cause is suspected
    \item \textbf{Reassurance and education:} confirming that the condition is real, common, and not a sign of mental illness
\end{itemize}

\section*{Management}
\begin{itemize}
    \item \textbf{Reassurance:} explaining that the hallucinations are caused by vision loss, are not a mental illness, and often improve over time
    \item \textbf{Improved lighting:} brightening the environment, particularly in the evening, can reduce the frequency of episodes
    \item \textbf{Stimulation:} keeping mentally and socially active may lessen episodes
    \item \textbf{Eye movements:} many people find that looking away, blinking, or moving the eyes makes an image disappear
    \item \textbf{Treating the eye condition:} improving vision where possible, such as through cataract surgery, may reduce symptoms
    \item \textbf{Low-vision support:} access to low-vision aids and rehabilitation services to make the most of remaining sight
    \item \textbf{Distraction:} turning attention to another activity, moving to a brighter room, or seeking company
    \item \textbf{Medication:} rarely, a specialist may consider medicines that reduce hallucinations if symptoms cause severe distress
    \item \textbf{Family education:} helping carers understand the condition so they respond calmly and supportively
\end{itemize}

\section*{Complications}
\begin{itemize}
    \item Distress, anxiety, and fear caused by the hallucinations
    \item Reluctance to leave the house or to be alone
    \item Social withdrawal and isolation
    \item Misdiagnosis as dementia or a psychiatric illness
    \item Sleep disturbance when hallucinations occur at night
    \item Falls if the person is distracted or startled while walking
    \item Reduced confidence and independence
    \item Reluctance to report symptoms, which delays reassurance and support
\end{itemize}

\section*{Prognosis and Outcomes}
Charles Bonnet syndrome is not dangerous, and the hallucinations do not indicate damage to the mind. In many people the episodes become less frequent or resolve within 12 to 18 months, although some people continue to experience them for many years. Symptoms often diminish when vision stabilises or improves, and they tend to be less troublesome once the person understands the cause and stops worrying about them. The syndrome itself does not progress to dementia or mental illness. The principal challenge is usually the anxiety it produces, which settles well with explanation, reassurance, and simple coping strategies.

\section*{Prevention and Self-Care}
\begin{itemize}
    \item Maintain good lighting throughout the home, especially in the evening
    \item Keep mentally and physically active to reduce episodes
    \item Attend regular eye checks and treat any reversible cause of vision loss
    \item Tell a doctor about the hallucinations so they can be properly explained
    \item Move the eyes, blink, or look away when an image appears
    \item Seek company or change activity if hallucinations feel intrusive
    \item Rest well, since tiredness can increase symptoms
    \item Use low-vision aids and support services to make the most of remaining sight
    \item Reassure family and carers that the condition is not a mental illness
\end{itemize}

\section*{Sources and Further Reading}
The following reputable sources provide further information on Charles Bonnet syndrome:
\begin{itemize}
    \item NHS. \textit{Health A--Z: conditions affecting vision and the eye.} https://www.nhs.uk/conditions/
    \item Mayo Clinic. \textit{Diseases and conditions library.} https://www.mayoclinic.org/
    \item MSD Manual. \textit{Professional edition: symptoms and causes of hallucinations.} https://www.msdmanuals.com/
    \item MedlinePlus. \textit{Health topics on vision and eye health.} https://medlineplus.gov/
    \item World Health Organization. \textit{Vision and eye health programmes.} https://www.who.int/
\end{itemize}
